\documentclass[11pt,a4paper]{article}
\usepackage[margin=2.3cm]{geometry}
% Built with tectonic (XeTeX): UTF-8 input and Latin Modern are the defaults.
\usepackage{microtype}
\usepackage{enumitem}
\usepackage{titlesec}
\usepackage[hidelinks]{hyperref}
\setlist{nosep,leftmargin=1.4em}
\setlength{\parskip}{0.45em}
\setlength{\parindent}{0pt}
\titleformat{\section}{\large\bfseries}{\thesection.}{0.5em}{}
\titlespacing*{\section}{0pt}{1.0em}{0.3em}

\newcommand{\studentname}{Furkan Aykut Sarıkamış}
\newcommand{\studentid}{26142007}

\begin{document}

\begin{center}
{\Large\bfseries Does Depth Confidence Help?\\[0.2em]
Confidence-Weighted Depth Supervision for Sparse-View 3D Gaussian Splatting}\\[0.8em]
Advanced Computer Vision --- Project Proposal\\[0.3em]
Group members: \studentname{} (\studentid), individual project
\end{center}

\section{Motivation and Problem Statement}

3D Gaussian Splatting (3DGS)~[1] reconstructs a scene from photographs and renders new viewpoints in real time. It works well with a hundred or more input images, but many real uses have only a handful: a robot that sees a place once, a quick inspection capture, a few phone photos of an object. With three input views the optimisation is under-constrained, and 3DGS overfits the training images while producing wrong geometry and floating artifacts in new views.

A common remedy is to supervise the rendered depth with a pretrained depth-estimation network~[3]. Recent studies report that this helps only marginally, and sometimes hurts, because the predicted depth is wrong in exactly the regions where guidance is needed most, such as thin structures, reflective or textureless surfaces and depth edges~[6]. Current depth foundation models such as Depth Anything~3~[4] also output a per-pixel confidence map. One recent method, Pi-GS~[5], weights its depth loss by such a confidence map, but does not measure what the weighting contributes on its own, and to my knowledge no study has checked whether these confidence maps actually indicate where the depth is wrong.

This project asks two specific questions: \emph{how much does depth supervision from a current depth foundation model improve a strong sparse-view 3DGS baseline, and does weighting that supervision by the model's own confidence add anything, given how reliably the confidence indicates depth error?}

\section{Objectives}

\begin{itemize}
\item Reproduce a strong published sparse-view 3DGS baseline (DropGaussian, CVPR 2025~[2]) on the LLFF 3-view benchmark.
\item Improve the baseline by adding depth supervision from a Vision-Transformer depth foundation model (Depth Anything~3).
\item Test whether weighting the depth loss by the model's confidence helps further, with a controlled ablation: no depth, unweighted depth, confidence-weighted depth, a binary confidence mask, and a shuffled-confidence control, each with three random seeds.
\item Evaluate with PSNR, SSIM and LPIPS on held-out views, and measure how well confidence predicts depth error using sparsification curves.
\item Analyse where the method helps and where it fails, and discuss applications.
\end{itemize}

\section{Methodology}

\textbf{Problem definition and dataset.} The task is novel view synthesis from three posed images. The main benchmark is LLFF~[7] (8 real forward-facing scenes) with the standard 3-view split at 1/8 resolution. Mip-NeRF~360~[8] with 12 input views is the second benchmark, to test whether the result holds on unbounded scenes. Both datasets are public. Preprocessing follows the baseline: camera poses and an initial point cloud from COLMAP, and depth and confidence maps computed once per training image.

\textbf{Model development.} The baseline is DropGaussian, which regularises 3DGS by randomly dropping Gaussians during training and uses no external prior. The proposed improvement adds one loss term: a Pearson-correlation depth loss between rendered depth and Depth Anything~3 depth, which is invariant to the unknown scale and shift of the predicted depth. The study then tests a confidence-weighted form of the same loss, in which each pixel is weighted by its normalised confidence; with uniform weights it reduces exactly to the unweighted loss, so the two are directly comparable. Implementation is in PyTorch. Each scene is optimised for 10{,}000 iterations on a single 8\,GB GPU.

\textbf{Evaluation and analysis.} All variants share the same code, initialisation, hyperparameters and seeds, so that differences are attributable to the loss alone. Results are reported as mean and standard deviation over three seeds with a paired significance test. Hyperparameters are chosen on two scenes and then frozen. For calibration, confidence is compared with depth error measured against a reference reconstruction trained on all available views. The report will state plainly whether the weighting helps, including if the result is negative. In a first single-seed pilot on LLFF, the depth loss raised the baseline from 20.54 to 20.93\,dB PSNR, while the confidence-weighted form was not distinguishable from the unweighted one.

\textbf{Demonstration.} Side-by-side rendered fly-through videos of the baseline and the proposed method on the test scenes, depth and confidence visualisations, and a reconstruction from three self-captured phone photos.

\section{Expected Outcomes}

\begin{itemize}
\item A report in IEEE proceedings format with the full quantitative evaluation, ablations and failure analysis.
\item Trained scene models for the baseline and each variant, with code and configurations to reproduce every table.
\item A short demo video and an interactive viewer of the reconstructed scenes.
\end{itemize}

\textbf{Timeline (3 weeks).} Week 1: environment, baseline reproduction, depth and confidence pipeline, loss implementation. Week 2: main experiments, second benchmark, calibration analysis. Week 3: demo, report, code clean-up.

\section*{References}
\small
\begin{enumerate}[label={[\arabic*]},leftmargin=2em]
\item B. Kerbl, G. Kopanas, T. Leimk\"uhler, G. Drettakis. 3D Gaussian Splatting for Real-Time Radiance Field Rendering. \emph{ACM Trans. Graphics (SIGGRAPH)}, 2023.
\item H. Park, G. Ryu, W. Kim. DropGaussian: Structural Regularization for Sparse-view Gaussian Splatting. \emph{CVPR}, 2025.
\item Z. Zhu, Z. Fan, Y. Jiang, Z. Wang. FSGS: Real-Time Few-shot View Synthesis using Gaussian Splatting. arXiv:2312.00451, 2023.
\item H. Lin, S. Chen, J. Liew, D. Y. Chen, Z. Li, G. Shi, J. Feng, B. Kang. Depth Anything 3: Recovering the Visual Space from Any Views. arXiv:2511.10647, 2025.
\item M. Hofer, M. Steinberger, T. K\"ohler. Pi-GS: Sparse-View Gaussian Splatting with Dense $\pi^3$ Initialization. arXiv:2602.03327, 2026.
\item W. Xiao, E. Goan, R. Santa Cruz, D. Ahmedt-Aristizabal, O. Salvado, C. Fookes, L. Lebrat. In Depth We Trust: Reliable Monocular Depth Supervision for Gaussian Splatting. arXiv:2604.05715, 2026.
\item B. Mildenhall, P. P. Srinivasan, R. Ortiz-Cayon, N. K. Kalantari, R. Ramamoorthi, R. Ng, A. Kar. Local Light Field Fusion: Practical View Synthesis with Prescriptive Sampling Guidelines. \emph{ACM Trans. Graphics (SIGGRAPH)}, 2019.
\item J. T. Barron, B. Mildenhall, D. Verbin, P. P. Srinivasan, P. Hedman. Mip-NeRF 360: Unbounded Anti-Aliased Neural Radiance Fields. \emph{CVPR}, 2022.
\end{enumerate}

\end{document}
